% Part: first-order-logic
% Chapter: syntax-and-semantics
% Section: main-operator

\documentclass[../../../include/open-logic-section]{subfiles}

\begin{document}

\olfileid{fol}{syn}{mai}

\olsection{सूत्राचा \printtoken{S}{main operator}}

\begin{explain}
!!a{formula}~$!A$ तयार करताना सर्वांत शेवटी वापरलेल्या कारकाविषयी
बोलणे अनेकदा उपयुक्त असते. या कारकाला $!A$ चा \emph{मुख्य कारक}
म्हणतात. अंतःप्रेरणेने तो $!A$ चा ``सर्वांत बाहेरचा'' कारक होय.
उदाहरणार्थ, $\lnot !A$ चा मुख्य कारक $\lnot$, $(!A \lor !B)$ चा
मुख्य कारक $\lor$, इत्यादी.
\end{explain}


\begin{defn}[!!^{main operator}]
\ollabel{def:main-op}
!!a{formula}~$!A$ चा \emph{!!{main operator}} पुढीलप्रमाणे परिभाषित केला जातो:
\begin{enumerate}
\item \indcase*{!A}{!A}{$\indfrm$ ला !!{main operator} नाही.}

\tagitem{prvNot}{\indcase{!A}{\lnot !B}{$\indfrm$ चा !!{main operator}
  हा~$\lnot$ आहे.}}{}

\tagitem{prvAnd}{\indcase{!A}{(!B \land !C)}{$\indfrm$ चा
  !!{main operator} हा~$\land$ आहे.}}{}

\tagitem{prvOr}{\indcase{!A}{(!B \lor !C)}{$\indfrm$ चा
  !!{main operator} हा~$\lor$ आहे.}}{}

\tagitem{prvIf}{\indcase{!A}{(!B \lif !C)}{$\indfrm$ चा
  !!{main operator} हा~$\lif$ आहे.}}{}

\tagitem{prvIff}{\indcase{!A}{(!B \liff !C)}{$\indfrm$ चा
  !!{main operator} हा~$\liff$ आहे.}}{}

\tagitem{prvAll}{\indcase{!A}{\lforall[x][!B]}{$\indfrm$ चा
  !!{main operator} हा~$\lforall$ आहे.}}{}

\tagitem{prvEx}{\indcase{!A}{\lexists[x][!B]}{$\indfrm$ चा
  !!{main operator} हा~$\lexists$ आहे.}}{}
\end{enumerate}
\end{defn}

प्रत्येक बाबतीत सूत्रातील !!{main operator} ची नेमकी दर्शवलेली
\emph{आस्थिति} आपल्याला अभिप्रेत आहे. उदाहरणार्थ,
$((!D \lif !E) \lif (!E \lif !D))$ हे सूत्र $(!B \lif !C)$ या रूपात
असून $!B$ हे $(!D \lif !E)$ आणि $!C$ हे $(!E \lif !D)$ असल्यामुळे,
$\lif$ ची दुसरी आस्थिति !!{main operator} आहे.

\begin{explain}
ही अशा फलनाची \emph{पुनरावर्ती} व्याख्या आहे जे सर्व अनाण्विक
!!{formula}s ना त्यांच्या !!{main operator} च्या आस्थितीकडे पाठवते.
!!{formula}s ची व्याख्या विगमनाने केलेली असल्यामुळे प्रत्येक
!!{formula}~$!A$ ला \olref{def:main-op} मधील एखादी बाब लागू होते.
त्यामुळे प्रत्येक अनाण्विक !!{formula}~$!A$ साठी !!a{main
  operator}
अस्तित्वात आहे याची खात्री मिळते. प्रत्येक !!{formula} ला यांपैकी
फक्त एकच अट लागू होते आणि प्रत्येक बाबतीत $!A$ ज्यांपासून तयार केले
आहे ती लहान !!{formula}s एकमेव रीतीने निर्धारित असतात. त्यामुळे $!A$
मधील !!{main
  operator} ची आस्थिति एकमेव असते आणि म्हणून आपण फलन
परिभाषित केले आहे.
\end{explain}

!!{formula}s चा !!{main operator} कोणते चिन्ह आहे यानुसार आपण त्यांना
\olref{tab:main-op} मधील नावे देतो.\iftag{defNot,defOr,defAnd,defIf,defIff,defTrue,defFalse,defEx,defAll}
{ मात्र लक्षात ठेवा की परिभाषित कारक औपचारिकदृष्ट्या !!{formula}s मध्ये
दिसत नाहीत. ते केवळ संक्षेप आहेत; म्हणून औपचारिकदृष्ट्या ते सूत्राचा
मुख्य कारक असू शकत नाहीत. त्यामुळे सर्व !!{formula}s विषयीच्या
सिद्धतांत त्यांचा स्वतंत्रपणे विचार करावा लागत नाही.}

\begin{table}[!h]
\centering
\begin{tabular}{c | c | c}
!!^{main operator} & !!{formula} चा प्रकार & उदाहरण\\
\hline
काही नाही & आण्विक (!!{formula}) &
\iftag{prvFalse}{$\lfalse$,}{}
\iftag{prvTrue}{$\ltrue$,}{}
$\Atom{R}{t_1, \dots, t_n}$,
$\eq[t_1][t_2]$\\
$\lnot$ & नकरण & $\lnot !A$ \\
$\land$ & संधियोग & $(!A \land !B$) \\
$\lor$ & विकल्पयोग & $(!A \lor !B$) \\
$\lif$ & !!{conditional} & $(!A \lif !B$) \\
$\liff$ & !!{biconditional} & $(!A \liff !B)$ \\
$\lforall[][]$ & वैश्विक (!!{formula})& $\lforall[x][!A]$ \\
$\lexists[][]$ & अस्तित्ववाची (!!{formula})& $\lexists[x][!A]$
\end{tabular}
\caption{!!{formula}s चे मुख्य कारक आणि नावे}
\ollabel{tab:main-op}
\end{table}

\end{document}
